\subsection{Forward prediction under geological shift}
\label{sec:tradeoff}
The mileage-ordered split evaluates the model on a later continuous tunnel
interval and therefore preserves the geological shift encountered during
excavation.  Under this setting, every model has negative $R^2$.  The proposed
framework obtains MAE 0.435 and is statistically indistinguishable from the
tested baselines.  It should therefore not be presented as a superior
forecasting model on the current dataset.

The ablations reinforce this cautious interpretation.  Removing monitoring
input degrades prediction, and removing the geometric prior produces the worst
learned-model result.  However, randomising the rock--machine edges improves
MAE from 0.435 to 0.413.  Thus the current experiment supports a contribution
from monitoring history and from the prior within the chosen architecture, but
does not demonstrate that the constrained edge topology improves forward
prediction.  The small test set and strong geological shift make this a
particularly demanding evaluation.

\subsection{Spatial evidence and its limits}
The saved-checkpoint diagnostics show positive spatial clustering on the TBM
surface (mean Moran's $I=0.529$).  This establishes that the projected relevance
is spatially structured rather than spatially random.  Component-level
differentiation is weak (mean CV\,$=0.018$), so the current evidence does not
support a strong claim that the model reliably separates cutterhead and shield
interaction roles.

Mean relevance is strongly inversely associated with TSP $V_p$ over the eight
future-test chainages (Pearson $r=-0.988$, $p<0.001$).  Although the direction is
consistent with stronger relevance in lower-velocity ground, this short
contiguous interval also contains strong chainage trends.  The correlation may
therefore reflect common spatial trend rather than a direct geological response.
It is retained as a local diagnostic, not as causal or physical validation.

The degree-control correlation is moderately negative ($r=-0.225$), indicating
that high relevance is not simply produced by more incident candidate edges.
At the same time, degree and relevance are not independent, and additional
topology-shuffle and spatial-trend controls are needed before stronger
interpretability claims are made.

\subsection{Application scenario: spatial reasoning beyond prediction}
\label{sec:application}
The main value currently supported by the framework is representational.  It
connects geological context, monitoring response, candidate rock--machine
relations, and TBM surface locations within one model and can project learned
relevance back to hotspot and chainage views.  Monitoring-only sequence models
cannot produce an equivalent TBM-surface diagnostic because they do not encode
surface nodes or rock--machine relations.

These views remain model diagnostics.  They are not calibrated risk estimates,
contact pressures, jamming probabilities, or prescriptions for machine control.
Operational use would require independent event labels, expert assessment, or
physical measurements that can validate whether the displayed hotspots
correspond to actual interaction states.

\subsection{Boundary and caution}
The case study contains only 44 effective sequence samples and 8 future-test
samples.  The resulting confidence intervals are broad, and the paired tests do
not establish a prediction advantage.  The test interval also differs
geologically from the training interval, so the negative $R^2$ values should be
read as evidence of limited extrapolation rather than ignored as an inconvenient
artifact.

The spatial evidence is checkpoint-based and post-hoc.  Multi-seed stability,
label permutation, candidate-edge threshold sensitivity, spatial detrending,
and cross-project validation remain necessary.  Learned relevance represents
response-consistent model importance under geometric screening, not a direct
measurement of contact force, stress, or engineering risk.

\subsection{Broader GIScience implications}
Despite these limitations, the framework demonstrates a GIScience-relevant
representation strategy for heterogeneous spatial entities under incomplete
observability.  Rock voxels, TBM surface components, candidate relations,
monitoring responses, and projected spatial diagnostics remain linked in one
explicit entity-relation model.

The results also show why spatial graph methods require two distinct evaluation
axes.  Prediction metrics test forward response generalisation, whereas spatial
diagnostics test whether learned relevance has coherent geographic structure.
The current framework provides spatially clustered diagnostics, but neither
predictive superiority nor physical validity is yet established.  Keeping these
claims separate is essential for a defensible interpretation.
